	\chapter{Scalar Products and Orthogonality}


	\section{Scalar Products}
	\begin{definitionenv}
		Suppose $V$ is a vector space over a field $\mathcal{K}(=\mathbb{R})$. A \textbf{scalar product} on $V$ is a mapping
		\[
		\langle \cdot, \cdot \rangle : V \times V \to \mathcal{K}
		\]
		such that for all $\mathbf{u}, \mathbf{v}, \mathbf{w} \in V$ and all $\alpha, \beta \in \mathcal{K}$, the following properties hold:
		\begin{enumerate}
			% \item (Conjugate symmetry) $\langle \mathbf{u}, \mathbf{v} \rangle = \overline{\langle \mathbf{v}, \mathbf{u} \rangle}$.
			\item (Symmetry) $\langle \mathbf{u}, \mathbf{v} \rangle = \langle \mathbf{v}, \mathbf{u} \rangle$.
			\item (Linearity in the first argument) $\langle \alpha \mathbf{u} + \beta \mathbf{v}, \mathbf{w} \rangle = \alpha \langle \mathbf{u}, \mathbf{w} \rangle + \beta \langle \mathbf{v}, \mathbf{w} \rangle$.
			% \item (Positive-definiteness) $\langle \mathbf{v}, \mathbf{v} \rangle \geq 0$, with equality if and only if $\mathbf{v} = \mathbf{0}$.
		\end{enumerate}
		% Then we call the pair $(V, \langle \cdot, \cdot \rangle)$ a \textbf{scalar product space}. 
		Additionally, we say the scalar product is \textbf{non-degenerate} if
		\[
		\langle \mathbf{v}, \mathbf{u} \rangle = 0, \quad \forall \mathbf{u} \in V \implies \mathbf{v} = \mathbf{0}.
		\]
	\end{definitionenv}
	\begin{remarkenv}
		The linearilty in the second argument also holds due to the conjugate symmetry.
	\end{remarkenv}
	\begin{definitionenv}
		Suppose $V$ is a vector space over a field $\mathcal{K}$ with a scalar product $\langle \cdot, \cdot \rangle$. We say two vectors $\mathbf{u}, \mathbf{v} \in V$ are \textbf{orthogonal} or \textbf{perpendicular} if
		\[
		\langle \mathbf{u}, \mathbf{v} \rangle = 0.
		\]
		Moreover, suppose $S \subseteq V$ is nonempty. We define the \textbf{orthogonal set} of $S$, as
		\[
		S^\perp = \{ \mathbf{v} \in V \mid \langle \mathbf{v}, \mathbf{s} \rangle = 0, \quad \forall \mathbf{s} \in S \}.
		\]
		And we say $\mathbf{v}\in S^\perp$ is \textbf{orthogonal} or \textbf{perpendicular} to $S$ and write $\mathbf{v} \perp S$.
	\end{definitionenv}
	\begin{propositionenv}
		Suppose $V$ is a vector space over a field $\mathcal{K}$ with a scalar product $\langle \cdot, \cdot \rangle$. If a set $S \le V$, then $S^\perp$ is a subspace of $V$. We often call $S^\perp$ the \textbf{orthogonal complement} of $S$ in $V$.
	\end{propositionenv}
	\begin{exampleenv}
		Suppose $A\in \mathbf{M}_{m \times n}(\mathcal{K})$, then $\mathbf{x}$ is a solution of the homogeneous linear system
		\[
		A\mathbf{x} = \mathbf{0}
		\]
		if and only if $\mathbf{x} \in (\operatorname{Row}(A))^\perp$, where $\operatorname{Row}(A)$ is the row space of matrix $A$. That is,
		\[
		\operatorname{Row}(A) \perp \operatorname{Null}(A).
		\]
	\end{exampleenv}
	\begin{definitionenv}
		Suppose $V$ is a finite-dimensional vector space over a field $\mathcal{K}$ with a scalar product $\langle \cdot, \cdot \rangle$. Then $V$ has a basis and any basis $\mathcal{B} = \{\mathbf{v}_1, \mathbf{v}_2, \dots, \mathbf{v}_n\}$ of $V$ is called an \textbf{orthogonal basis} if
		\[
		\langle \mathbf{v}_i, \mathbf{v}_j \rangle = 0, \quad \forall i \neq j\in \{1,2,\ldots,n\}.
		\]

	\end{definitionenv}
	\begin{definitionenv}
		Suppose $V$ is a vector space over a field $\mathcal{K\left( =\mathbb{R} \right)}$ with a scalar product $\langle \cdot, \cdot \rangle$. We call this scalar product \textbf{positive-definite} if
		\[
		\langle \mathbf{v}, \mathbf{v} \rangle 
		\begin{cases}
		> 0, & \mathbf{v} \neq \mathbf{0}, \\
		= 0, & \mathbf{v} = \mathbf{0}.
		\end{cases}
		\]
		If the scalar product is positive-definite, we define the \textbf{norm} of a vector $\mathbf{v} \in V$ as
		\[
		\| \mathbf{v} \| = \sqrt{\langle \mathbf{v}, \mathbf{v} \rangle}.
		\]
		Then we can define the \textbf{distance} between two vectors $\mathbf{u}, \mathbf{v} \in V$ as
		\[
		dist(\mathbf{u}, \mathbf{v}) = \| \mathbf{u} - \mathbf{v} \|.
		\]
		Also, if $\mathbf{v} \in V$ satisfies $\| \mathbf{v} \| = 1$, we say $\mathbf{v}$ is a \textbf{unit vector}. Note that for any nonzero vector $\mathbf{v} \in V$, the vector
		\[
		\frac{\mathbf{v}}{\| \mathbf{v} \|}
		\]
		is a unit vector called the \textbf{normalization} of $\mathbf{v}$.
	\end{definitionenv}
	\begin{remarkenv}
		Here we write $\mathcal{K} \left(= \mathbb{R}\right)$ because we need ``$>0$'' to define positive-definiteness. Actually $\mathcal{K}$ can be any field homomorphic to $\mathbb{R}$. 
	\end{remarkenv}
	\begin{exampleenv}
		Let $f,g$ be real-valued functions integrable on the interval $I$ containing $[0,1]$. We define
		\[
		\langle f, g \rangle = \int_0 ^1 f(x) g(x) \,\mathrm{d}x, \quad \left\| f \right\| = \sqrt{\langle f, f \rangle}.
		\]
		Then $\langle \cdot, \cdot \rangle$ is a positive-definite scalar product on the vector space of all real-valued functions integrable on the interval $I$.
	\end{exampleenv}
	\begin{propositionenv}
		Suppose $V$ is a vector space over a field $\mathcal{K}\left( = \mathbb{R} \right)$ with a positive-definite scalar product $\langle \cdot, \cdot \rangle$.
		\begin{enumerate}
			\item For $c\in \mathcal{K}$ and $\mathbf{v} \in V$, we have $\| c \mathbf{v} \| = |c| \, \| \mathbf{v} \|$.
			\item For any $\mathbf{u}, \mathbf{v} \in V$,
			\[
			\| \mathbf{u}-\mathbf{v} \| = \| \mathbf{u}+\mathbf{v} \| \quad \text{if and only if} \quad \mathbf{u} \perp \mathbf{v}.
			\]
			\item (Pythagorean theorem) For any $\mathbf{u}, \mathbf{v} \in V$ with $\mathbf{u} \perp \mathbf{v}$,
			\[
			\| \mathbf{u} + \mathbf{v} \|^2 = \| \mathbf{u} \|^2 + \| \mathbf{v} \|^2.
			\]
			\item (Parallelogram law) For any $\mathbf{u}, \mathbf{v} \in V$,
			\[
			\| \mathbf{u} + \mathbf{v} \|^2 + \| \mathbf{u} - \mathbf{v} \|^2 = 2 \left( \| \mathbf{u} \|^2 + \| \mathbf{v} \|^2 \right).
			\]
			\item (Cauchy–Schwarz inequality) For any $\mathbf{u}, \mathbf{v} \in V$,
			\[
			|\langle \mathbf{u}, \mathbf{v} \rangle| \leq \| \mathbf{u} \| \, \| \mathbf{v} \|.
			\]
			\item (Triangle inequality) For any $\mathbf{u}, \mathbf{v} \in V$,
			\[
			\| \mathbf{u} + \mathbf{v} \| \leq \| \mathbf{u} \| + \| \mathbf{v} \|.
			\]
			\item Suppose $\mathbf{u}, \mathbf{v} \in V$ and $\mathbf{u} \neq \mathbf{0}$. Then
			\[
			c = \frac{\langle \mathbf{v}, \mathbf{u} \rangle}{\langle \mathbf{u}, \mathbf{u} \rangle}
			\]
			is called the \textbf{component} of $\mathbf{v}$ in the direction of $\mathbf{u}$ and we call $c \mathbf{u}$ the \textbf{projection} of $\mathbf{v}$ onto $\mathbf{u}$. Note that 
			\[
			\left( \mathbf{v} - c \mathbf{u} \right) \perp \mathbf{u}.
			\]
			Furthermore,
			\[
			\| \mathbf{v} - c \mathbf{u} \| = \min_{\alpha \in \mathcal{K}} \| \mathbf{v} - \alpha \mathbf{u} \|.
			\]
			\item More generally, suppose $\mathbf{u_1}, \mathbf{u_2}, \dots, \mathbf{u_k} \in V$ are non-zero vectors that are pairwise orthogonal. For any $\mathbf{v} \in V$, we define the \textbf{component} of $\mathbf{v}$ along $u_i$ to be 
			\[
			c_i = \frac{\langle \mathbf{v}, \mathbf{u_i} \rangle}{\langle \mathbf{u_i}, \mathbf{u_i} \rangle}, \quad i=1,2,\ldots,k.
			\]
			We call $c_i \mathbf{u_i}$ the \textbf{projection} of $\mathbf{v}$ onto $\mathbf{u_i}$. Note
			\[
			\langle \mathbf{v} - \sum_{i=1}^k c_i \mathbf{u_i}, \mathbf{u_j} \rangle = \langle \mathbf{v}, \mathbf{u_j} \rangle - \sum_{i=1}^k c_i \langle \mathbf{u_i}, \mathbf{u_j} \rangle = 0, \quad j=1,2,\ldots,k.
			\]
			and we have
			\[
			\left( \mathbf{v}- \sum_{i=1}^k c_i \mathbf{u_i} \right) \perp \mathbf{u}, \quad \forall \mathbf{u}\in \operatorname{span}\left\{ \mathbf{u_1}, \mathbf{u_2}, \ldots, \mathbf{u_k} \right\}.
			\]
			We also call $\sum\limits_{i=1}^k c_i \mathbf{u_i}$ the \textbf{projection} of $\mathbf{v}$ onto $\operatorname{span}\left\{ \mathbf{u_1}, \mathbf{u_2}, \ldots, \mathbf{u_k} \right\}$.
			Furthermore,
			\[
			\left\| \mathbf{v} - \sum_{i=1}^k c_i \mathbf{u_i} \right\| = \min_{\alpha_1, \alpha_2, \ldots, \alpha_k \in \mathcal{K}} \left\| \mathbf{v} - \sum_{i=1}^k \alpha_i \mathbf{u_i} \right\|.
			\]
		\end{enumerate}
	\end{propositionenv}
	\begin{definitionenv}
		Suppose $V$ is a finite-dimensional vector space over a field $\mathcal{K}\left( = \mathbb{R} \right)$ with a scalar product $\langle \cdot, \cdot \rangle$. A basis $\mathcal{B} = \{\mathbf{v}_1, \mathbf{v}_2, \dots, \mathbf{v}_n\}$ of $V$ is called an \textbf{orthonormal basis} if
		\[
		\langle \mathbf{v}_i, \mathbf{v}_j \rangle =
		\begin{cases}
		1, & i = j, \\
		0, & i \neq j,
		\end{cases}
		\quad \forall i,j \in \{1,2,\ldots,n\}.
		\]
	\end{definitionenv}
	\begin{theoremenv}
		Suppose $V$ is a finite-dimensional vector space over a field $\mathcal{K}\left( = \mathbb{R} \right)$ with a positive-definite scalar product $\langle \cdot, \cdot \rangle$. Let $\{w_1, w_2, \ldots, w_k\}$ be an orthogonal set of vectors in $V$. Then for any $\mathbf{v} \in V$, we have
		\[
		\langle \mathbf{v}, \mathbf{v} \rangle \geq \sum_{i=1}^k \frac{|\langle \mathbf{v}, w_i \rangle|^2}{\langle w_i, w_i \rangle}.
		\]
		with equality if and only if $\mathbf{v} \in \operatorname{span}\{w_1, w_2, \ldots, w_k\}$.

		This is called \textbf{Bessel's inequality}.
	\end{theoremenv}
	\begin{lemmaenv}
		Suppose $V$ is a finite-dimensional vector space over $\mathcal{K}\left( = \mathbb{R} \right)$ with a positive-definite scalar product $\langle \cdot, \cdot \rangle$. Let $W$ be a proper subspace of $V$. If $\{w_1, w_2, \ldots, w_k\}$ is an orthonormal basis of $W$, then there exist vectors $w_{k+1}, w_{k+2}, \ldots, w_n \in V$ such that $\{w_1, w_2, \ldots, w_k, w_{k+1}, w_{k+2}, \ldots, w_n\}$ is an orthonormal basis of $V$.
	\end{lemmaenv}
	\begin{proofenv}
		This is a constructive proof via the Gram–Schmidt process. 

		Note that we do know there exists $x_{k+1}, x_{k+2}, \ldots, x_n \in V$ such that $\{w_1, w_2, \ldots, w_k, x_{k+1}, x_{k+2}, \ldots, x_n\}$ is a basis of $V$. Next we will convert this basis to an orthonormal basis via projection.

		\textbf{Step 1:} Project $x_{k+1}$ onto $W$ and denote the projection as
		\[
		w_{k+1} = x_{k+1} - \sum_{i=1}^k \frac{\langle x_{k+1}, w_i \rangle}{\langle w_i, w_i \rangle} w_i.
		\]
		Note that since
		\[
		\langle w_{k+1}, w_j \rangle = \langle x_{k+1}, w_j \rangle - \sum_{i=1}^k \frac{\langle x_{k+1}, w_i \rangle}{\langle w_i, w_i \rangle} \langle w_i, w_j \rangle = 0, \quad j=1,2,\ldots,k,
		\]
		we have $w_{k+1} \perp W$. Also, $w_{k+1} \neq \mathbf{0}$ since otherwise $x_{k+1} \in W$ which contradicts the fact that $\{w_1, w_2, \ldots, w_k, x_{k+1}, x_{k+2}, \ldots, x_n\}$ is a basis of $V$. Then we normalize $w_{k+1}$.

		\textbf{Inductive Step:} We proceed by induction to construct the following vectors $w_{k+2}, w_{k+3}, \ldots, w_n$.

	\end{proofenv}
	\begin{theoremenv}
		Suppose $V$ is a finite-dimensional nontrivial vector space over a field $\mathcal{K}\left( = \mathbb{R} \right)$ with a positive-definite scalar product $\langle \cdot, \cdot \rangle$. Then $V$ has an orthonormal basis.
	\end{theoremenv}
	\begin{theoremenv}
		Suppose $V$ is a vector space over a field $\mathcal{K}\left( = \mathbb{R} \right)$ with a positive-definite scalar product $\langle \cdot, \cdot \rangle$ and $\dim V = n$. Define $W$ be a nontrivial subspace of $V$ with $\dim W = r \left( 1\le r \le n-1 \right) $. Define the \textbf{orthogonal complement} of $W$ in $V$, which is a subspace of $V$, as
		\[
		W^\perp = \{ \mathbf{v} \in V \mid \langle \mathbf{v}, \mathbf{w} \rangle = 0, \quad \forall \mathbf{w} \in W \}.
		\]
		Then we have
		\[
		\dim W + \dim W^\perp = \dim V, \quad W \oplus W^\perp = V.
		\]
	\end{theoremenv}
    \section{Hermitian Products}
	% \begin{remarkenv}
		For vector space $V$ over $\mathbb{C}$, if we want to have a ``scalar product'' that is ``postive-definite'' so that
		\[
		\langle \mathbf{v}, \mathbf{v} \rangle \in \mathbb{R} \text{ and } \langle \mathbf{v}, \mathbf{v} \rangle \ge 0,
		\]
		we need to modify our definition of ``scalar product''.
	% \end{remarkenv}
	\begin{definitionenv}
		Suppose $V$ is a vector space over $\mathbb{C}$. A \textbf{Hermitian scalar product} on $V$ is a mapping
		\[
		\langle \cdot, \cdot \rangle : V \times V \to \mathbb{C}
		\]
		satisfying the following properties:
		\begin{enumerate}
			\item (Conjugate symmetry) $\langle \mathbf{u}, \mathbf{v} \rangle = \overline{\langle \mathbf{v}, \mathbf{u} \rangle}$.
			\item (Linearity in the first argument) $\langle \alpha \mathbf{u} + \beta \mathbf{v}, \mathbf{w} \rangle = \alpha \langle \mathbf{u}, \mathbf{w} \rangle + \beta \langle \mathbf{v}, \mathbf{w} \rangle$.
		\end{enumerate}
		Additionally, we say the Hermitian scalar product is \textbf{positive-definite} if
		\[
		\langle \mathbf{v}, \mathbf{v} \rangle 
		\begin{cases}
		> 0, & \mathbf{v} \neq \mathbf{0}, \\
		= 0, & \mathbf{v} = \mathbf{0}.
		\end{cases}
		\]
		Note that
		\[
		\langle \mathbf{v}, \mathbf{v} \rangle = \overline{\langle \mathbf{v}, \mathbf{v} \rangle} \implies \langle \mathbf{v}, \mathbf{v} \rangle \in \mathbb{R}.
		\]
		the positive-definiteness is well-defined.
	\end{definitionenv}
	\begin{remarkenv}
		Note that the linearity in the second argument does not hold. Instead, we have
		\[
		\langle \mathbf{w}, \alpha \mathbf{u} + \beta \mathbf{v} \rangle = \overline{\alpha} \langle \mathbf{w}, \mathbf{u} \rangle + \overline{\beta} \langle \mathbf{w}, \mathbf{v} \rangle.
		\]
		This is called the \textbf{conjugate linearity} in the second argument.
	\end{remarkenv}
	\begin{remarkenv}
		For vector space $V$ over $\mathbb{C}$, given a Hermitian product, we say $x,y$ are orthogonal if $\langle x,y \rangle = 0\left( =\langle y,x \rangle \right)$. Concepts like orthogonal complement can be defined accordingly. If additionally the Hermitian product is positive-definite, concepts like norm, distance, orthogonal basis, orthonormal basis, etc. can be defined accordingly.
	\end{remarkenv}
	\begin{definitionenv}
		Given $A\in \mathbf{M}_{m \times n}(\mathbb{R})$, we call $\operatorname{Im}(A)$ the \textbf{column space} of $A$ and $\operatorname{Im}(A^T)$ the \textbf{row space} of $A$. We also call $\dim \operatorname{Im}(A)$ the \textbf{column rank} of $A$ and $\dim \operatorname{Im}(A^T)$ the \textbf{row rank} of $A$. Since the column rank and the row rank are equal, we simply call it the \textbf{rank} of $A$ and denote it as $\operatorname{rank}(A)$.
	\end{definitionenv}
	\begin{remarkenv}
		Let $A\in \mathbf{M}_{m \times n}(\mathbb{R})$ and $R=\operatorname{rref}(A)$. Then we know $\ker(A) = \ker(R)$ and thus $\operatorname{rank}(A) = \operatorname{rank}(R)=\text{pivots of } R$.
	\end{remarkenv}

    % \newpage
    \section{Bilinearity and Quadratic Form}
    \begin{definitionenv}
        Suppose $U,V,W$ are vector spaces over a field $\mathcal{K}$ and $g: U \times V \to W$ be a map. We say $g$ is \textbf{bilinear} if for each fixed $u \in U$, the map $v \mapsto g(u,v)$ is linear from $V$ to $W$, and for each fixed $v \in V$, the map $u \mapsto g(u,v)$ is linear from $U$ to $W$.
    \end{definitionenv}
    \begin{exampleenv}
        Let $A \in \mathbf{M}_{m \times n}(\mathcal{K})$. Define $g_A: \mathcal{K}^m \times \mathcal{K}^n \to \mathcal{K}$ as $g_A(x,y) = x^T A y$. Then $g_A$ is a bilinear map.
    \end{exampleenv}
    \begin{theoremenv}
        Given any bilinear map $g: \mathcal{K}^m \times \mathcal{K}^n \to \mathcal{K}$, there exists a unique matrix $A \in \mathbf{M}_{m \times n}(\mathcal{K})$ such that $g(x,y) = x^T A y \left(=g_A\right)$ for all $x\in \mathcal{K}^m$ and $y\in \mathcal{K}^n$.
    \end{theoremenv}
    \begin{theoremenv}
        The set of all bilinear maps from $\mathcal{K}^m \times \mathcal{K}^n$ to $\mathcal{K}$, denoted by $\operatorname{Bil}(\mathcal{K}^m \times \mathcal{K}^n, \mathcal{K})$, is a vector space over $\mathcal{K}$. The association $A \mapsto g_A$ is a vector space isomorphism from $\mathbf{M}_{m \times n}(\mathcal{K})$ to $\operatorname{Bil}(\mathcal{K}^m \times \mathcal{K}^n, \mathcal{K})$. That is, we have
        \[
        \dim \operatorname{Bil}(\mathcal{K}^m \times \mathcal{K}^n, \mathcal{K}) = \dim \mathbf{M}_{m \times n}(\mathcal{K}) = mn.
        \]
    \end{theoremenv}
    \begin{remarkenv}
        Note that a scalar product define on a vector space $V$ is a bilinear map from $V \times V$ to $\mathcal{K}$. 
        
        Particularly, if $V = \mathcal{K}^n$, then the unique matrix $A\in \M_{n \times n}(\mathcal{K})$ corresponding to the scalar product is symmetric, i.e., $A^T = A$.
    \end{remarkenv}
    \begin{remarkenv}
        In some textbook, a scalar product is also called a symmetric bilinear form.
    \end{remarkenv}
    \begin{definitionenv}
        Suppose $V$ is a vector space over a field $\mathcal{K}$ and $\langle \cdot, \cdot \rangle$ is a scalar product on $V$. We call the map $f: V \to \mathcal{K}$ defined by $f(\mathbf{v}) = \langle \mathbf{v}, \mathbf{v} \rangle$ a \textbf{quadratic form} on $V$.
    \end{definitionenv}
    \begin{exampleenv}
        For the case $V=\mathcal{K}^n$, a quadratic form on $V$ can be uniquely associated with a symmetric matrix $A\in \mathbf{M}_{n \times n}(\mathcal{K})$ such that
        \[
        f(x) = x^T A x = \sum_{i=1}^{n}\sum_{j=1}^{n} x_i a_{ij} x_j \quad \forall x \in \mathcal{K}^n.
        \]
        Particularly, if $A$ is a diagonal matrix, then we have
        \[
        f(x) = \sum_{i=1}^{n} a_{ii} x_i^2.
        \]
    \end{exampleenv}
    
    \section{Determinants}
    Recall in last semester we defined the determinant of a $2\times 2$ matrix $\displaystyle A = \begin{bmatrix} a & b \\ c & d \end{bmatrix}$ as $\displaystyle \det(A) = \begin{vmatrix} a & b \\ c & d \end{vmatrix} = ad - bc$. And we conclude that $A$ is invertible if and only if $\det(A) \neq 0$. Furthermore, we know $\displaystyle A^{-1} = \frac{1}{\det(A)} \begin{bmatrix} d & -b \\ -c & a \end{bmatrix}$ when $A$ is invertible.

    We recall more properties of determinants of $2\times 2$ matrices:
    \begin{enumerate}
        \item $\det(I_2) = 1$.
        \item Bilinearity: For any $a,b,c,d \in \mathcal{K}$, we have
        \[
        \begin{vmatrix} a+b & c \\ d & e \end{vmatrix} = \begin{vmatrix} a & c \\ d & e \end{vmatrix} + \begin{vmatrix} b & c \\ d & e \end{vmatrix}, \quad \begin{vmatrix} a & b+c \\ d & e \end{vmatrix} = \begin{vmatrix} a & b \\ d & e \end{vmatrix} + \begin{vmatrix} a & c \\ d & e \end{vmatrix}.
        \]
        \item Alternating property: For any $a,b,c,d \in \mathcal{K}$, we have
        \(
        \begin{vmatrix} a & b \\ c & d \end{vmatrix} = -\begin{vmatrix} c & d \\ a & b \end{vmatrix}.
        \)
        \item $\det(A)=\det(A^T)$.
    \end{enumerate}
    Now we try to expand the definition of determinant to more general cases.
    \begin{definitionenv}
        Based on the definition of $2\times 2$ determinant, we define $3\times 3$ determinant as follows: suppose
        \[
        A = \begin{bmatrix}
        a_{11} & a_{12} & a_{13} \\
        a_{21} & a_{22} & a_{23} \\
        a_{31} & a_{32} & a_{33}
        \end{bmatrix} \in \mathbf{M}_{3 \times 3}(\mathcal{K}),
        \]
        then we define
        \[
        \det(A) = a_{11} \begin{vmatrix} a_{22} & a_{23} \\ a_{32} & a_{33} \end{vmatrix} - a_{12} \begin{vmatrix} a_{21} & a_{23} \\ a_{31} & a_{33} \end{vmatrix} + a_{13} \begin{vmatrix} a_{21} & a_{22} \\ a_{31} & a_{32} \end{vmatrix}.
        \]
    \end{definitionenv}
    \begin{remarkenv}
        We call this formula the \textbf{Laplace expansion} or the \textbf{cofactor formula} by the first row.
    \end{remarkenv}
    \begin{theoremenv}
        The $3\times 3$ determinant satisfies the following properties:
        \begin{enumerate}
            \item $\det(I_3) = 1$.
            \item Trilinearity: as a function of the columns in $A\in \M_{3 \times 3}(\mathcal{K})$, the determinant is linear in each column when the other two columns are fixed. For example, $\det(A_1, \alpha x + \beta y, A_3) = \alpha \det(A_1, x, A_3) + \beta \det(A_1, y, A_3)$ for any $A_1, A_3 \in \mathcal{K}^3$, $x,y \in \mathcal{K}^3$ and $\alpha, \beta \in \mathcal{K}$.
            \item If two adjacent columns of $A$ are equal, then $\det(A) = 0$.
        \end{enumerate}
    \end{theoremenv}
    \begin{definitionenv}
        Suppose function $D: \mathbf{M}_{n \times n}(\mathcal{K}) \to \mathcal{K}$ or alternatively $D: \mathcal{K}^n \times \mathcal{K}^n \times \cdots \times \mathcal{K}^n \to \mathcal{K}$ (with $n$ copies of $\mathcal{K}^n$) with notation $D(A)$ or $D(A_1, A_2, \ldots, A_n)$ where $A_i$ is the $i$-th column of $A$. We say:
        \begin{enumerate}
            \item $D$ is multilinear: as a function of the columns in $A\in \M_{n \times n}(\mathcal{K})$, $D$ is linear in each column when the other $n-1$ columns are fixed.
            \item $D$ is alternating: if two adjacent columns of $A$ are equal, then $D(A) = 0$.
        \end{enumerate}
    \end{definitionenv}
    \begin{theoremenv}
        There exists a unique multilinear and alternating function $D: \mathbf{M}_{n \times n}(\mathcal{K}) \to \mathcal{K}$ such that $D(I_n) = 1$. We call this unique $D$ the $n\times n$ \textbf{determinant} function and denote it as $\det(\cdot)$.
    \end{theoremenv}
    \begin{proofenv}
        We prove in a constructive and inductive way.

        \noindent
        \textbf{Existence:} 
        
        \textbf{Initial Cases}: For $n=1$, we define $\det([a]) = a$ for any $a \in \mathcal{K}$. For $n=2$ and $n=3$, we have already defined the determinant.

        \textbf{Inductive assumption}: suppose such $D$ exists as a map from $\mathbf{M}_{k \times k}(\mathcal{K})$ to $\mathcal{K}$ such that the three properties hold for any $ k \le n-1$.

        \textbf{Inductive step}: For
        \[
        A = \begin{bmatrix}
            a_{11} & \cdots & a_{1j} & \cdots & a_{1n} \\
            \vdots & \ddots & \vdots & \ddots & \vdots \\
            a_{i1} & \cdots & a_{ij} & \cdots & a_{in} \\
            \vdots & \ddots & \vdots & \ddots & \vdots \\
            a_{n1} & \cdots & a_{nj} & \cdots & a_{nn}
        \end{bmatrix} \in \mathbf{M}_{n \times n}(\mathcal{K}),
        \]
        we define
        \[
        D(A) = \sum_{j=1}^n (-1)^{i+j} a_{ij} D(\tilde{A}_{ij}),
        \]
        where $\tilde{A}_{ij}$ is the $(n-1)\times (n-1)$ matrix obtained by deleting the $i$-th row and the $j$-th column of $A$. By the inductive assumption, $D(\tilde{A}_{ij})$ is well-defined and satisfies the three properties. We will show that such $D$ satisfies the three properties as well.

        For multilinearity, we consider the individual term $(-1)^{i+j} a_{ij} D(\tilde{A}_{ij})$ in the sum and a particular $k\in \{1,2,\ldots,n\}$.
        \begin{itemize}
            \item If $j\neq k$, then $(-1)^{i+j} a_{ij} $ is independent of $A_k$ and $D(\tilde{A}_{ij})$ is linear in $A_k$ by the inductive assumption. Thus $(-1)^{i+j} a_{ij} D(\tilde{A}_{ij})$ is linear in $A_k$.
            \item If $j=k$, then $D(\tilde{A}_{ij})$ is independent of $A_k$, but $(-1)^{i+j} a_{ij}$ is linear in $A_k$. Thus $(-1)^{i+j} a_{ij} D(\tilde{A}_{ij})$ is linear in $A_k$.
        \end{itemize}
        Hence $D(A)$ is linear in $A_k$.

        For alternating property, suppose $A_k = A_{k+1}$ for some $k\in \{1,2,\ldots,n-1\}$. Consider:
        \begin{itemize}
            \item If $j\neq k$ and $j\neq k+1$, then $D(\tilde{A}_{ij})$ has two adjacent columns equal and thus $D(\tilde{A}_{ij})=0$ by the inductive assumption. Thus $(-1)^{i+j} a_{ij} D(\tilde{A}_{ij}) = 0$.
            \item In the other cases, we only have
            \[
            D(A) = (-1)^{i+k} a_{ik} D(\tilde{A}_{ik}) + (-1)^{i+k+1} a_{i,k+1} D(\tilde{A}_{i,k+1}) = 0.
            \]
        \end{itemize}

        For identity, we have
        \[
        D(I_n) = (-1)^{i+i} D(I_{n-1}) = 1.
        \]
        \textbf{Uniqueness:} This part is finished in Homework. The core idea is to utilize the permutation of columns and the alternating property to show that $D$ is determined by its value on $I_n$.
    \end{proofenv}
    \begin{theoremenv}
        Here are some additional properties of the determinant function:
        \begin{enumerate}
            \item If the $i$-th and $j$-th $(i\neq j)$ columns of $A$ are swapped, then $\det(A)$ changes its sign.
            \item If $A$ has two identical columns or rows, then $\det(A) = 0$.
            \item If one adds a scalar multiple of one column to another column, then $\det(A)$ does not change.
            \item $\det(A) = \det(A^T)$.
            \item $\det(AB) = \det(A)\det(B)$ for any $A,B \in \mathbf{M}_{n \times n}(\mathcal{K})$.
        \end{enumerate}
    \end{theoremenv}
    \begin{theoremenv}
        Suppose $A\in\M_{n \times n}(\mathcal{K})$ satisfies $\det(A) \neq 0$ and $b\in \mathcal{K}^n$. If $x\in \mathcal{K}^n$ is a solution to the linear system $Ax=b$, then $x$ satisfies
        \[
        x_i = \frac{\det(A_{ib})}{\det(A)}, \quad i=1,2,\ldots,n,
        \]
        where $A_{ib}$ is the matrix obtained by replacing the $i$-th column of $A$ with $b$. This is called \textbf{Cramer's rule}.
    \end{theoremenv}
    \begin{proofenv}
        Note $b = x_1 A_1 + x_2 A_2 + \cdots + x_n A_n$, where $A_i$ is the $i$-th column of $A$.
        Then we have
        \[
        \det(A_{ib}) = \det(A_1, \ldots, A_{i-1}, \sum_{j=1}^n x_j A_j, A_{i+1}, \ldots, A_n) = x_i \det(A).
        \]
    \end{proofenv}
    \begin{theoremenv}
        Suppose $A\in \mathbf{M}_{n \times n}(\mathcal{K})$. If $\det(A) \neq 0$, then the columns of $A$ are linearly independent and thus $A$ is invertible.
    \end{theoremenv}
    \begin{proofenv}
        It's equivalent to show that if the columns of $A$ are linearly dependent, then $\det(A) = 0$. Suppose there exists $k\in \{1,2,\ldots,n\}$ such that 
        \[
        A_k = \sum_{j=1, j\neq k}^n \alpha_j A_j,
        \]
        where $\alpha_j \in \mathcal{K}$ for $j=1,2,\ldots,n$ and $\alpha_j$ are not all zero. Then we have
        \[
        \det(A) = \det(A_1, \ldots, A_{k-1}, \sum_{j=1, j\neq k}^n \alpha_j A_j, A_{k+1}, \ldots, A_n) = 0.
        \]
    \end{proofenv}
    \begin{definitionenv}
        Consider the following two types of column operations:
        \begin{enumerate}
            \item Add a scalar multiple of one column to another column.
            \item Interchange two columns.
        \end{enumerate}
        We call two matrices $A,B \in \M_{n \times n}(\mathcal{K})$ are \textbf{column equivalent} if $B$ can be obtained from $A$ by a succession of column operations of the above two types.
    \end{definitionenv}
    \begin{theoremenv}
            Let $A,B \in \M_{n \times n}(\mathcal{K})$ be column equivalent. Then:
            \begin{enumerate}
                \item $\operatorname{rank} (A) = \operatorname{rank}(B)$.
                \item $A$ is invertible if and only if $B$ is invertible.
                \item $\det A = 0$ if and only if $\det B = 0$.
            \end{enumerate}
    \end{theoremenv}
    \begin{proofenv}
        For (1), we only need to show that the column operations of the above two types do not change the rank of a matrix. This is straightforward since these two types of column operations do not change the linear dependence or independence of columns.

        For (2), since $A$ is invertible, we have $\operatorname{rank}(A) = n$. By (1), we have $\operatorname{rank}(B) = n$ and thus $B$ is invertible.

        For (3), since the column operations only changes the sign of the determinant, we have $\det A = 0$ if and only if $\det B = 0$.
    \end{proofenv}
    \begin{theoremenv}
        Suppose $A\in \M_{n \times n}(\mathcal{K})$ is column equivalent to a lower-triangular matrix $L$.
    \end{theoremenv}
    \begin{theoremenv}
        Suppose $A\in \M_{n \times n}(\mathcal{K})$. The following conditions are equivalent:
        \begin{enumerate}
            \item $A$ is invertible.
            \item The columns of $A$ are linearly independent.
            \item $\det A \neq 0$.
        \end{enumerate}
    \end{theoremenv}
    \begin{remarkenv}
        To calculate $\det A$, we can:
        \begin{enumerate}
            \item Laplace expansion by a row or a column.
            \item Permutation definition.
            \item Convert $A$ to a lower-triangular or upper-triangular matrix $L$ by column operations of the above two types and then calculate $\det L$ as the product of diagonal entries of $L$.
        \end{enumerate}
    \end{remarkenv}

